% ------------------------------------------------------------------
%  Capa analítica: emplazamiento y dimensionamiento
%  Fragmento del Subdocumento 4.2, traido desde contenido.tex con
%  \parte{partes/13_l_capa_analitica}. Las rutas son relativas a la carpeta del
%  subdocumento, nunca a la de main.tex.
% ------------------------------------------------------------------

\chapter{Capa analítica: emplazamiento y dimensionamiento}
\label{sec:l-capa-analitica-emplazamiento-y-dimension}

Complemento de la sección de dimensionamiento y plan de capacidad --- Emplazamiento y dimensionamiento de la capa analítica. La capa analítica no es una funcionalidad añadida por iniciativa de la propuesta: es un componente exigido por el numeral 5.4 de las Transversales (capa analítica separada de la transaccional, tableros operacionales y de gestión con filtros por período y unidad organizacional, desagregación hasta el hecho y latencia máxima declarada) y por el Capítulo 18 del caso, que fija OTIF, fill rate y costo de servir como criterios de aceptación. Esta sección declara con qué componentes se materializa, dónde se emplazan y cómo se garantiza que ninguna consulta analítica degrade la operación.

Alcance de lo que no se propone. Conforme a la decisión de diseño sobre analítica avanzada del registro de decisiones, la solución \textbf{no incorpora inteligencia artificial ni analítica predictiva} en el alcance contratado, y lo declara de forma fundada como permite el Capítulo 18 de las Transversales. La capa analítica se entrega preparada ---lago de datos en formato columnar y almacén analítico--- para incorporarlas cuando el CLIENTE alcance la madurez de datos que hoy no tiene: el caso parte de un 41 \% de recepciones sin registro de lote y un 2,3 \% de diferencia de inventario, y predecir sobre esa base produce confianza injustificada. El ruteo se resuelve con optimización determinística, auditable y corregible por el planificador, que es lo que pide el criterio de aceptación N\textdegree{} 7 del caso.

\section{Fundamento normativo}
\label{sub:2-fundamento-normativo}

\subsection{Exigencias de las Bases Técnicas Transversales}

\begin{itemize}
  \item \textbf{RT-05.25} --- Capa analítica con tableros operacionales y de gestión (obligatorio; estructura de plataforma).
  \item \textbf{RT-05.26} --- Tableros con filtros por período, unidad organizacional y drill-down (obligatorio; estructura de plataforma).
  \item \textbf{RT-14.02} --- Acceso propio y permanente a tableros operacionales y de negocio (obligatorio; observabilidad).
  \item \textbf{RT-14.03} --- SLA/SLO medidos sobre experiencia real del usuario (P95) (obligatorio; observabilidad).
  \item \textbf{RT-14.04} --- Alertamiento basado en síntomas de negocio (obligatorio; observabilidad).
\end{itemize}

El numeral 5.4 de las Transversales hace de la capa analítica un componente obligatorio y no una funcionalidad adicional: RT-05.25 a RT-05.28 son obligatorios y RT-05.29 fija la latencia máxima según el caso. La analítica predictiva de RT-05.30 es deseable y esta propuesta la declina de forma fundada, conforme a RT-18.01.

\subsection{Exigencias del Caso 02}

El caso define tres indicadores estratégicos que requieren procesamiento analítico continuo:

\begin{itemize}
  \item \textbf{OTIF} (entregas completas y a tiempo): estado actual 82,4 \%, meta $>$ 95 \% (Cap. 4.1 y 7.1).
  \item \textbf{Fill rate}: estado actual 91,3 \%, meta $>$ 97 \% (Cap. 4.1).
  \item \textbf{Costo de servir por cliente}: actualmente desconocido, meta calcular y gestionar (Cap. 7.2 y 9.8).
\end{itemize}

Estos indicadores no pueden calcularse sin un componente de analítica que consolide datos transaccionales de múltiples módulos (preventa, transporte, cobranza, telemetría, bodega) y los presente de forma consumible por los gerentes.

\section{Requerimientos que dependen de la capa analítica}
\label{sub:3-requerimientos-que-dependen-de-la-capa-a}

En total 25 requerimientos alimentan este módulo, todos fichados en el catálogo del Subdocumento 3: 10 propios de la Épica 11 (RF-11.01 a RF-11.08 + RNF-11.01 latencia $\leq$ 5 min + RNF-11.02 desacople OLAP/OLTP), 4 transversales (RF-16.02 tableros del CLIENTE, RF-16.03 alertas de negocio, RF-17.12 exportación, RNF-16.01 SLA P95) y 11 de otras épicas que producen o consumen datos analíticos (costos de ruta, cierre de turno, rendición, excursión térmica, envases, kilometraje, ocupación de zona y alerta de crédito).

\section{Mapeo de requerimientos a componentes de arquitectura}
\label{sub:4-mapeo-de-requerimientos-a-componentes-de}

\subsection{Componentes de la capa analítica}

El módulo de BI/Analítica se materializa en los componentes que se listan en la Tabla~\ref{tab:t95}:

\begin{tablalafrox}{Componentes del módulo BI}{tab:t95}%
  {F{0.24} F{0.23} F{0.24} Y}%
  {\cab{Componente} & \cab{Ubicación} & \cab{Tecnología} & \cab{Responsabilidad}}
  Motor OLAP / Warehouse & Cloud (AWS) & Amazon Redshift Serverless & Almacenamiento analítico desacoplado de OLTP. Consolida datos de Aurora PostgreSQL vía DMS CDC. Alimenta RF-11.01--RF-11.08 \\
  Motor transaccional fuente & Cloud (AWS) & Amazon Aurora PostgreSQL (PostGIS) & Base OLTP de donde se extraen los datos hacia Redshift. Contiene datos de preventa, reparto, cobranza, bodega \\
  Motor transaccional fuente (on-prem) & On-Premise (Talca) & PostgreSQL + PostGIS (local) & WMS de bodega. Sincroniza con Aurora vía DMS CDC (RPO $\leq$ 15 min) \\
  Capa de procesamiento analítico & Cloud (AWS) & Celery workers (Django) + Redshift SQL & Cálculo de OTIF (RF-11.03), Fill Rate (RF-11.04), Costo de Servir (RF-11.02), Ocupación de Flota (RF-11.05), Segmentación (RF-11.08) \\
  Distribución de reportes & Cloud (AWS) & Celery Beat + SES (correo) + S3 (almacenamiento PDF/XLSX) & RF-11.01: reportes automáticos semanales/mensuales \\
  Tablero operacional & Cloud (AWS) & Grafana (autoservicio) o Angular SPA & RF-11.06: tablero de control en tiempo real. RF-16.02: acceso del cliente. RF-16.03: alertas de negocio \\
  Sistema de alertas & Cloud (AWS) & CloudWatch + SNS + Celery & RF-16.03: alertamiento por síntomas de negocio. RF-11.05: alertas de ocupación \\
  Portal de autoatención cliente & Cloud (AWS) & Angular SPA + API Gateway & RF-16.02: tableros para clientes con datos en tiempo real y exportación \\
  Caché de indicadores & Cloud (AWS) & Amazon ElastiCache (Redis) & RNF-11.01: latencia $\leq$ 5 min para despliegue de indicadores \\
\end{tablalafrox}

\subsection{Flujo de datos analíticos}

En texto, el flujo de datos analíticos es: fuentes transaccionales (Preventa App, Bodega WMS, Transporte GPS y Cobranza App) $\rightarrow$ motor transaccional Aurora PostgreSQL (Preventa, Inventario, Transporte y Cobranza, cada dominio OLTP) $\rightarrow$ DMS CDC (extracción continua e incremental) $\rightarrow$ motor analítico Redshift Serverless (dimensiones Tiempo, Cliente, Ruta, Producto, Vehículo y Zona; hechos fact\_entregas, fact\_costos y fact\_inventario; vistas materializadas mv\_otif\_diario, mv\_costo\_servir, mv\_fill\_rate, mv\_ocupacion\_flota y mv\_rentabilidad\_cliente) $\rightarrow$ dos salidas: Tablero Gerencial (Grafana/Angular --- RF-11.06 operacional, RF-16.02 para el cliente, RF-16.03 alertas) y Distribución por correo (SES + Celery Beat --- informes semanal y mensual, RF-11.01).

\subsection{Desacople OLAP/OLTP}

Redshift Serverless (OLAP) opera independiente de Aurora PostgreSQL (OLTP); DMS CDC extrae cambios de forma incremental y continua sin bloquear transacciones. Los indicadores se calculan como vistas materializadas en Redshift (refresco cada 5--15 min) y se cachean en ElastiCache Redis para tableros de alta concurrencia. Resultado: picking mantiene $\leq$ 1s y preventa $\leq$ 1.5s mientras los gerentes consultan tableros simultáneamente.

\section{Componente de costo de servir}
\label{sub:5-componente-de-costo-de-servir}

El cálculo del Costo de Servir (RF-11.02) es el requerimiento analítico más complejo del módulo, ya que consolida datos de 4 fuentes distintas, que se desglosan en la Tabla~\ref{tab:t96}:

\begin{tablalafrox}{Componente de Costo de Servir}{tab:t96}%
  {F{0.24} F{0.24} F{0.24} Y}%
  {\cab{Componente de costo} & \cab{Fuente de datos} & \cab{Módulo productor} & \cab{Fórmula simplificada}}
  Combustible devengado & Telemetría GPS (litros/km) & Telemetría (Épica 14) & km recorridos $\times$ precio litro $\times$ rendimiento vehicular \\
  Peajes & Registro de viaje & Transporte (Épica 6) & Monto real de peajes por viaje \\
  Tiempo de conducción & Telemetría GPS (horas) & Telemetría (Épica 14) & Horas de conducción $\times$ costo hora conductor \\
  Tiempo de descarga & App móvil de reparto & Transporte (Épica 6) & Minutos de descarga $\times$ costo hora \\
  Costo de bodega (prorrateo) & WMS & Almacén (Épica 2) & Costo total bodega $\div$ volumen almacenado $\times$ volumen cliente \\
  Costo administrativo (prorrateo) & Facturación ERP & Administración & Costo fijo administrativo $\div$ N pedidos $\times$ pedidos cliente \\
  Costo de devoluciones/mermas & Logística inversa & Logística Inversa (Épica 8) & Costo de envases perdidos + productos dañados \\
  Costo de mermas por vencimiento & Inventario & Almacén (Épica 2) & Valor de producto vencido imputado al cliente \\
\end{tablalafrox}

La fórmula de rentabilidad neta resultante (RF-11.08):

Margen neto = Ingreso por venta $-$ Costo de mercadería $-$ Costo de servir real

Donde el Costo de servir real = $\Sigma$(componentes anteriores) por cliente y por entrega.

\section{Componente de alertas de negocio}
\label{sub:6-componente-de-alertas-de-negocio}

Las alertas por síntomas de negocio se distinguen de las alertas de infraestructura en que miden impacto al negocio, no estado de servidores, como se detalla en la Tabla~\ref{tab:t97}:

\begin{tablalafrox}{Componente de Alertas de Negocio (RF-16.03)}{tab:t97}%
  {F{0.24} F{0.24} Y F{0.24}}%
  {\cab{Síntoma de negocio} & \cab{Umbral configurable} & \cab{Fuente de datos} & \cab{Acción}}
  OTIF diario $<$ meta & $<$ 95\% (configurable) & Redshift (mv\_otif\_diario) & Notificación gerente comercial + calidad \\
  Fill Rate mensual $<$ meta & $<$ 97\% (configurable) & Redshift (mv\_fill\_rate) & Alerta gerente comercial \\
  Costo de servir $>$ umbral & $>$ \$X por entrega (configurable) & Redshift (mv\_costo\_servir) & Alerta gerente finanzas \\
  Ocupación de flota $<$ umbral & $<$ 65\% (configurable) & Redshift (mv\_ocupacion\_flota) & Alerta planificador de rutas \\
  Excursión térmica en tránsito & $>$ umbral por tipo producto & IoT/DynamoDB + SNS (Lambda validador) & Alerta conductor + calidad \\
  Crédito de cliente excedido & Deuda + pedido $>$ cupo & Aurora PostgreSQL (OLTP) & Bloqueo en app de preventa \\
  Envases retornables excedidos & Saldo $>$ umbral configurable & Aurora PostgreSQL (OLTP) & Alerta preventista \\
\end{tablalafrox}

\section{Requerimientos de tableros}
\label{sub:7-requerimientos-de-tableros-rt-05-25-rt-0}

Las Bases exigen tableros con las capacidades que se listan en la Tabla~\ref{tab:t98}:

\begin{tablalafrox}{Requerimientos de tableros}{tab:t98}%
  {F{0.26} G{0.11} Y}%
  {\cab{Capacidad exigida} & \cab{RT} & \cab{Implementación}}
  Tableros operacionales y de gestión & RT-05.25 & Dashboards separados: operacional (tiempo real, RF-11.06) y estratégico (periódico, RF-11.01) \\
  Filtros por período & RT-05.26 & Filtros de rango de fechas en Grafana/Angular (hoy, semana, mes, trimestre, año) \\
  Filtros por unidad organizacional & RT-05.26 & Filtros por zona, ruta, sucursal, cliente, canal de comercialización \\
  Drill-down & RT-05.26 & Navegación de resumen $\rightarrow$ detalle: total $\rightarrow$ zona $\rightarrow$ ruta $\rightarrow$ cliente $\rightarrow$ entrega individual \\
  Datos en tiempo real & RT-14.02 & Actualización cada $\leq$ 5 min (RNF-11.01) vía Redis cache + Redshift refresh \\
  Acceso del cliente & RT-14.02 & Portal web Angular con autenticación Keycloak, RF-16.02 \\
  Exportación & RT-16.28 & Descarga en PDF, XLSX, CSV con filtros aplicados, RF-17.12 \\
\end{tablalafrox}

\subsection{Tablero operacional (RF-11.06)}

El tablero operacional cubre los paneles que se listan en la Tabla~\ref{tab:t99}:

\begin{tablalafrox}{Tablero operacional (RF-11.06)}{tab:t99}%
  {F{0.24} Y F{0.24} F{0.24}}%
  {\cab{Panel} & \cab{Indicadores mostrados} & \cab{Fuente} & \cab{Frecuencia actualización}}
  Avance de rutas & Pedidos completados / total, \% avance por ruta & Aurora PostgreSQL (OLTP) & Tiempo real ($<$ 1 min) \\
  Entregas cumplidas & \% OTIF del día, entregas a tiempo vs atrasadas & Redshift (mv\_otif\_diario) & $\leq$ 5 min \\
  Devoluciones en tránsito & Cantidad de devoluciones, productos devueltos & Aurora PostgreSQL (OLTP) & Tiempo real \\
  Alertas cadena de frío & Excursiones térmicas activas, historial 24h & DynamoDB (IoT) + SNS (Lambda validador) & Tiempo real \\
  Recaudación del día & Monto recaudado en efectivo, cheques, transferencias & Aurora PostgreSQL (cobranza) & Tiempo real \\
\end{tablalafrox}

\subsection{Tablero gerencial (RF-11.01)}

Distribución automática de dos reportes ejecutivos: semanal (OTIF, Fill Rate, Costo de Servir por zona, devoluciones) y mensual (los anteriores más merma por vencimiento, rentabilidad por cliente y desglose por canal). Destinatarios configurables; distribución por correo vía SES + Celery Beat (PDF/XLSX).

\section{Cumplimiento de Bases Administrativas}
\label{sub:8-cumplimiento-de-bases-administrativas}

\begin{itemize}
  \item \textbf{Separación transaccional/analítico} (Art. 23\textdegree{}). Almacén analítico alimentado por extracción incremental, sin consulta alguna contra el transaccional en vivo (subsección de justificación del desacople entre lo analítico y lo transaccional).
  \item \textbf{Exportabilidad total} (Art. 23\textdegree{} \textperiodcentered{} RT-05.06). Todo informe y tablero es exportable en formatos abiertos, con el modelo semántico documentado para que el CLIENTE construya sus propios informes (RT-05.27 y RT-05.28).
  \item \textbf{Tableros para el CLIENTE} (Art. 25\textdegree{} \textperiodcentered{} RT-14.02). Tres niveles de tablero ---operacional, gerencial y del mandante---; el del mandante en solo lectura y con auditoría de consultas (subsección de requerimientos de tableros).
  \item \textbf{Alertamiento por síntomas de negocio} (Art. 25\textdegree{} \textperiodcentered{} RT-14.03). Componente de alertas de negocio sobre OTIF, fill rate, excursión térmica y descuadre de rendición (subsección de componente de alertas de negocio).
  \item \textbf{Gestión de capacidad} (Art. 25\textdegree{} \textperiodcentered{} RT-09.09). Revisión trimestral del consumo real contra el proyectado, con aviso al 70 \% de ocupación (la sección de dimensionamiento y plan de capacidad, subsección de actualización del plan de capacidad).
  \item \textbf{Analítica predictiva} (RT-05.30, deseable). No se propone en el alcance contratado, con declaración fundada conforme a RT-18.01; la capa queda preparada para incorporarla cuando el CLIENTE alcance la madurez de datos.
\end{itemize}

\section{Cobertura y cierre}
\label{sub:9-cobertura-y-cierre}

En total, 25 requerimientos alimentan al módulo BI: 10 propios de la Épica 11 (8 RF + 2 RNF), 4 transversales (3 RF + 1 RNF) y 11 RF de otras épicas. Todos son trazables a su origen normativo.
